% \iffalse meta-comment
%
% Copyright (C) 2020 by Martin C Schwarzer <polyluxus@gmail.com>
% ---------------------------------------------------------------------------
% This work may be distributed and/or modified under the
% conditions of the LaTeX Project Public License, either version 1.3
% of this license or (at your option) any later version.
% The latest version of this license is in
%   http://www.latex-project.org/lppl.txt
% and version 1.3 or later is part of all distributions of LaTeX
% version 2005/12/01 or later.
%
% This work has the LPPL maintenance status `maintained'.
%
% The Current Maintainer of this work is Martin C Schwarzer.
%
% This work consists of the files polycv.dtx and polycv.ins
% and the derived filebase polycv.cls.
%
% The code is available at <https://github.com/polyluxus/polycv>.
%
% \fi
%
% \iffalse
%<*driver>
\ProvidesFile{polycv.dtx}
%</driver>
%<class>\NeedsTeXFormat{LaTeX2e}[1999/12/01]
%<class>\ProvidesClass{polycv}
%<*class>
    [2020/06/29 v1.0rc1 Curriculum Vitae Class]
%</class>
%
%<*driver>
\documentclass{ltxdoc}
\PassOptionsToPackage{dvipsnames}{xcolor}
\usepackage{dtxdescribe}
\definecolor{lightgray}{gray}{0.95}
\usepackage[T1]{fontenc}
\usepackage[utf8]{inputenc}
\usepackage{libertine}
\usepackage[english]{babel}
\usepackage{csquotes}
\usepackage{setspace}
\usepackage{listings}
\lstdefinestyle{mytex}{
  language=[LaTeX]TeX,%
  backgroundcolor=\color{lightgray},%
  basicstyle=\footnotesize\ttfamily,%
  keywordstyle=\bfseries,%
  keywordsprefix=polycv,%
  breaklines=true,%
  morekeywords={%
    author, title, position, email, phone, mobile, github, orcid, address, location, street,
    insertauthor, inserttitle, insertposition, insertemail, insertphone, insertmobile, 
    insertgithub, insertorcid, insertaddress, insertlocation, insertstreet,
    signaturefile, insertsigfilename, orcidiconfilename,
    cvline, cvitemline, cventry, cvskill, cvlanguage, cvlineaddress, cvlineemail, 
    cvlinemobile, cvlinephone, cvlineorcid, cvlinegithub,
    progressbar, levelbarsquares, levelbarcircles,
    setmargins, setheaderheight, setfooterheight, seticonspace, sethintcolwidth,
    sethintcolsep, setcolprimary, setcolsecondary, setcolhighlight, setcolshade,
    setprogbarheight, setprogbarunits, setheaderfont, setentryragged, setentryjustified,
    setsigleft, setsigright,
    opening, closing, ps,
    pdfappendix,
    section, subsection,
    includegraphics, faTerminal, faStar
  },%
  escapeinside={(*@}{@*)}%
}
\lstset{style=mytex}
\usepackage{fontawesome}
\usepackage[%
  pdftex, bookmarks=true, hidelinks,%
  colorlinks, linkcolor=MidnightBlue, urlcolor=MidnightBlue,%
  pageanchor=true, hyperindex=true,
]{hyperref}
\hypersetup{%
  pdfinfo={%
    Title={The LaTeX class PolyCV},%
    Author={Martin C Schwarzer},%
    Subject={A Curriculum Vitae Class},%
    keywords={LaTeX, CV, Letter, Curriculum Vitae}}}

\newcommand\mybackslash{\char`\\}
\newcommand\myarg[1]{\(\langle\)\textit{#1}\(\rangle\)}

\setcounter{IndexColumns}{2}

\EnableCrossrefs % \DisableCrossrefs %
\CodelineIndex
\RecordChanges
\begin{document}
  \DocInput{polycv.dtx}
  \PrintChanges
  \PrintIndex
\end{document}
%</driver>
% \fi
%
% \CheckSum{0}
%
% \CharacterTable
%  {Upper-case    \A\B\C\D\E\F\G\H\I\J\K\L\M\N\O\P\Q\R\S\T\U\V\W\X\Y\Z
%   Lower-case    \a\b\c\d\e\f\g\h\i\j\k\l\m\n\o\p\q\r\s\t\u\v\w\x\y\z
%   Digits        \0\1\2\3\4\5\6\7\8\9
%   Exclamation   \!     Double quote  \"     Hash (number) \#
%   Dollar        \$     Percent       \%     Ampersand     \&
%   Acute accent  \'     Left paren    \(     Right paren   \)
%   Asterisk      \*     Plus          \+     Comma         \,
%   Minus         \-     Point         \.     Solidus       \/
%   Colon         \:     Semicolon     \;     Less than     \<
%   Equals        \=     Greater than  \>     Question mark \?
%   Commercial at \@     Left bracket  \[     Backslash     \\
%   Right bracket \]     Circumflex    \^     Underscore    \_
%   Grave accent  \`     Left brace    \{     Vertical bar  \|
%   Right brace   \}     Tilde         \~}
%
%
% \changes{v1.0}{2020/06/18}{Initial version}
%
% \GetFileInfo{polycv.dtx}
%
% \DoNotIndex{\providecommand,\newcommand,\renewcommand,\newenvironment,\def,\edef,\let,\par}
% \DoNotIndex{\AfterPreamble, \AtEndPreamble, \ProcessKeyvalOptions, \SetupKeyvalOptions}
% \DoNotIndex{\IfFileExists,}
% \DoNotIndex{\#,\$,\%,\&,\@,\\,\{,\},\^,\_,\~,\ ,\newline,\noindent,\strut}
% \DoNotIndex{\@ne, \@author, \@title, \today}
% \DoNotIndex{\begin,\end,\bgroup,\egroup,\begingroup,\endgroup,\newcount,\advance}
% \DoNotIndex{\catcode,\closein,\closeout,\if,\ifnum,\else,\fi,\empty,\detokenize}
% \DoNotIndex{\baselineskip,\linewidth,\color,\colorlet,\hfill,\href,\hspace,\vspace,%
%   \loop,\repeat,\relax,\rule}
% \DoNotIndex{\ignorespaces,\makebox,\newlength,\setlength,\parbox,\parindent,%
%   \raisebox,\resizebox}
% \DoNotIndex{\number,\numexpr}
% \DoNotIndex{\includegraphics,\includepdf,\pagecolor,\nopagecolor,\pagestyle,\overfullrule}
% \DoNotIndex{\geometry,\newgeometry,\restoregeometry,\paperwidth}
% \DoNotIndex{\bfseries,\sffamily,\raggedright,\raggedleft,\flushright,\centering,%
%   \small,\footnotesize,\Huge,\Large,\large,\normalsize,\textbf,\textit}
% \DoNotIndex{\ClassError,\ClassWarning,\ClassInfo,\LoadClass,\PassOptionsToClass,%
%   \PassOptionsToPackage,\requirePackage}
% \DoNotIndex{\CurrentOptionKey,\DeclareBoolOption,\DeclareComplementaryOption,%
%   \DeclareDefaultOption,\DeclareFieldFormat,\DeclareStringOption,\DeclareVoidOption}
% \DoNotIndex{\draw,\filldraw,\node}
% \DoNotIndex{\biblabelsep}
%
% \title{The \pkg{polycv} class}
% \author{Martin C Schwarzer\\\href{mailto:polyluxus@gmail.com}{polyluxus@gmail.com}}
% \date{\fileversion~from \filedate}
%
% \maketitle
%
% \begin{abstract}
%   \pkg{Polycv} is a document class to typeset curriculum vitaes with \LaTeX.
%   It also provides a template for a cover letter consistent with the layout
%   of the resume.
% \end{abstract}
%
% \tableofcontents
% \clearpage
%
% \section{Introduction}
%
% I previously used the package
% \href{https://ctan.org/pkg/koma-moderncvclassic}{koma-moderncvclassic}
% to prepare my CV, and I was actually quite satisfied with the look of
% it. In a way this class borrows some design elements from it. However, I
% noticed it was not using a lot of space, which got especially
% problematic when inserting a picture.
% 
% I looked through various templates and classes online (%
% \href{https://github.com/opieters/limecv}{limecv},
% \href{https://github.com/latex-ninja/simple-academic-resume}{simple academic resume},
% \href{https://github.com/deedy/Deedy-Resume}{DeedyResume},
% \href{https://www.sharelatex.com/templates/cv-or-resume/fancy-cv}{FancyCV by Adrien Friggeri},
% \href{https://www.overleaf.com/gallery/tagged/cv}{the Overleaf CVGallery}, 
% TeX Stack Exchange \href{https://tex.stackexchange.com/q/80/33413}{%
%   LaTeX template for resume/curriculum vitae}, \ldots{} , there are possibly some I forgot),
% but couldn't find anything that was exactly to my linking. So I decided
% to write a class for myself, not just only for my CV, but also for the
% training in writing such classes (or packages).
% 
% The main goal was to structure the information in a way that, if
% condensed enough, it would fit on one page; but additional pages could
% be added when the need arose, e.g.~a long publication list, or many
% stations along the way.
% 
% Additionally I liked the idea of having a cover letter template, which
% resembles most of the elements of the CV.
% 
% The source code is available at \href{https://github.com/polyluxus/polycv}{\faGithub~GitHub},
% where there are also a few example uses of this class.
% If anything is unclear, needs improvement, is missing please submit an issue there
% (\href{https://github.com/polyluxus/polycv/issues}{GitHub issues}). Thank you.
%
% \subsection{Contributing}
%
% If you find this class useful, you are very welcome to contribute to it.
% The repository is hosted at 
% \href{https://github.com/polyluxus/polycv}{\faGithub~GitHub repository}.
% It will consist of one main branch, the release branch, possibly a release candidate
% branch and a developement branch.
% If you would like to contribute changes, please submit a pull request to
% the development branch. Once all changes have been pulled in, the commits will be squashed 
% and pulled into the release candidate branch for final testing and adjustment.
% As such, there will always be only one commit for one release.
% 
% \section{Installation}
% 
% This work consists of the files \lstinline`polycv.dtx` and \lstinline`polycv.ins`
% and the derived filebase \lstinline`polycv.cls`, which should be copied to a location
% where latex will look for it.
% How to build the package and documentation from its source files is explained in
% section~\ref{sec:building}.
% The source code is hosted at \href{https://github.com/polyluxus/polycv}{\faGithub~GitHub},
% where the derived class and documentation files are also available.
% 
% \subsection{Prerequisites}
% In order for this class to function properly, the following packages must be installed:
% \begin{itemize}
%   \item The package \lstinline`etoolbox` provides additional hooks necessary for
%     the setup within the source class.
%   \item The \lstinline`calc` package allows to derive some lengths from others.
%   \item The \lstinline`hyperref` package provides clickable cross-references,
%     e.g. to references to the bibliography, or websites.
%   \item Colorsupport comesss from \lstinline`xcolor`.
%   \item Most of the graphical elements within this class are produced with 
%     \lstinline`tikz`, such as the header and footer boxes, and the progressbar.
%   \item To include a signature and profile pictures, \lstinline`graphicx` is lloaded.
%   \item For most of the document the indentation is disabled with the 
%     \lstinline`parskip` package.
%   \item The page layout is produced with the \lstinline`geometry` package.
%   \item The \lstinline`fontawesome` is needed for icons, e.g. \faGithub, in the side bar.
%   \item To include other reference letters as pdfs, \lstinline`pdfpages` is also loaded.
% \end{itemize}
%
% \clearpage
% \section{Usage}
%
% You can load the class in the same way you would load any other class.
% The following listing defines a skeleton document containing a elements of the class.
% \iffalse
%<*example>
% \fi
\begin{lstlisting}
\documentclass[(*@\myarg{options}@*)]{polycv}
  %
  % packages, 
  % define personal information,
  % define colours, page layout, etc.
  %
\begin{document}
  \begin{polycvfirstpage}
    \begin{polycvsidebar}
      %
      % Contents of the side bar box
      % (personal information image, etc.)
      %
    \end{polycvsidebar}
    %
    % Body of the first page
    %
  \end{polycvfirstpage}
  \begin{polycvpage}
    %
    % Body of additional pages
    %
  \end{polycvpage}
  \begin{polycvletter}%
    [(*@\myarg{Author Address}@*)]%
    {(*@\myarg{Subject}@*)}%
    {(*@\myarg{Recipient Address}@*)}
    \opening{Dear Example,}
    %
    % Letter content
    %
    \closing{Sincerely,}
    \ps{more text}
  \end{polycvletter}
  \pdfappendix{(*@\myarg{filename}@*)}
\end{document}
\end{lstlisting}
% \iffalse
%</example>
% \fi
% \subsection{Class Options}
% This class is build on top of the standard article class. 
% It therefore inherits most options from that class. 
% Some are incompatible with the design of this class and have been disabled.
% Other options allow for customisation of the used templates.
% In this section they will be briefly described, along with commands that can be 
% used to set the values in the preamble.
% 
% \subsubsection{Page Layout}\label{sec:page-layout}
%
% The general page layout will be accomplished by the \lstinline`geopetry` package. 
% Depending on the environment currently in use, a few internal lengths are defined and used.
% Most of these can only be set via the class options, or within the preamble.
% 
% \DescribeOption{margins}
% \DescribeDefault{1.0 cm}
% The length of \lstinline`margins` is the baselength of most surrounding space,
% it defines and sets the length \lstinline`\polycvmargins`. 
% It is used additionally to the header and footer lengths.
% This setting also controls the space between the side bar and the main body on the 
% first page of the CV environment.
% \DescribeMacro{\setmargins} 
% The command \lstinline`setmargins`\marg{length} is only available in the preable 
% and can be used to set the margin length after loading the class.
% 
% \DescribeOption{headerheight}
% \DescribeDefault{3.0cm}
% \DescribeOption{footerheight}
% \DescribeDefault{0.5cm}
% The options \lstinline`headerheight` and \lstinline`footerheight` control the height 
% of the boxes for the header and the footer, respectively.
% These lengths affect the whole document and can therfore only be set when loading the
% class or within the preamble.
% The will be the same for each of the environments used.
% \DescribeMacro{\setheaderheight}
% \DescribeMacro{\setfooterheight}
% Instead of choosing the key/value option at loadtime, the commands
% \lstinline`\setheaderheight`\marg{length} and 
% \lstinline`\setfooterheight`\marg{length} can be issued to manipulate these values.
%
% \DescribeOption{sidebarwidth}
% \DescribeDefault{5.0cm}
% On the left side of first page of the CV a box is created, the side bar, which may hold 
% personal information, a picture, skills, languages, hobbies, etc..
% The width of this box is controlled with the option \lstinline`sidebarwidth`.
% This must be defined before the sidebar is attempted to be set.
% \watchout
% \DescribeMacro{\setsidebarwidth}
% There is currently no command analogous to the other lengths to set this in the preable,
% it must be defined via the options.
% 
% \DescribeOption{iconspace}
% \DescribeDefault{0.8cm}
% Closely related is the option \lstinline`iconspace`, which defines the width that is 
% reserved for the icon in a line in the side bar environment.
% \DescribeMacro{\seticonspace}
% It can also be set via \lstinline`\seticonspace`\marg{length} in the preamble.
%
% \DescribeOption{hintcol}
% \DescribeDefault{2.7cm}
% \DescribeOption{hintcolsep}
% \DescribeDefault{0.2cm}
% In the main body, each entry will consist of two parts: 
% First, a small column for a keyword, date, bar, which is called hint in this documentation.
% The width of this column can be adjusted with the option \lstinline`hintcol`.
% Second, a column for the description of this hint. This is just a parbox, and depends on 
% the  command used. The width will be automaticall calculated from the available space.
% These two columns are separated by a space.
% The width of which can be set via the option \lstinline`hintcolsep`.
% \DescribeMacro{\sethintcolwidth}
% \DescribeMacro{\sethintcolsep}
% Alternatively they can be set in the preamble via 
% \lstinline`\sethintcolwidth`\marg{length} and
% via \lstinline`\sethintcolsep`\marg{length}, respectively.
%
% \DescribeOption{letterindent} 
% \DescribeDefault{1em}
% In the letter environment the indentation of the first line of a paragraph can be adjusted
% with the \lstinline`letterindent` option.
% This is another length that can (currently) only be controlled with passing the option 
% to the class.
% 
% \subsubsection{Colors}
% \DescribeOption{xcolor}\DescribeDefault{dvipsnames}
% The class loads the \lstinline`xcolor` package and passes by default the 
% option \lstinline`dvipsnames` to it.
% This can be changed by using the \lstinline`xcolor` option to pass other parameters 
% to this package.
% It is loaded internally before the \lstinline`tikz` package to avoid option conflicts.
% 
% \DescribeOption{primary}\DescribeDefault{blue}
% The main colour of the CV is called \lstinline`primary` and can be set via this option.
% It is used for all parts that will be highlighted, such as the header bar and the 
% section headings.
% \DescribeOption{shade} \DescribeDefault{black}
% \DescribeOption{mixing} \DescribeDefault{10}
% The textcolor is a composit color from the primary color and the \lstinline`shade` color. 
% By default this is defined as black, so the text will appear darker.
% The amount of primary color mixed in depends on the a factor which can be set via 
% the \lstinline`mixing` option.
% 
% \DescribeOption{secondary} \DescribeDefault{white}
% The \lstinline`secondary` color specifies the background color and is white by default.
%
% \margintag{unused}
% \DescribeOption{highlight} \DescribeDefault{orange} 
% One last color is defined to highlight items. 
% It can be set via the \lstinline`highlight` option.
% This option currently has little effect as the color itself is not used within the template.
% 
% \DescribeMacro{\setcolprimary}
% \DescribeMacro{\setcolsecondary}
% \DescribeMacro{\setcolhighlight}
% \DescribeMacro{\setcolshade}
% To allow for custom colors to be used, all colors ca also be set with the 
% \lstinline`\setcol<...>`\marg{color} commands.
% Note that the mixing parameter must be set via the options and has no special macro.
%
% \subsubsection{Progress and Level Bar Options}
% 
% \DescribeOption{progbarheight}
% \DescribeDefault{1.5 ex}
% Skills can be displayd as an CV entry and accompanied with a progress bar or a
% level (bar) indicator. The width of this bar is by default either determined by the
% environment in which it is used, or the width of the hint column.
% The height of the bar can be controlled with the \lstinline`progbarheight` option.
% \DescribeMacro{\setprogbarheight}
% It can also be set via \lstinline`\setprogbarheight`\marg{length}.
% \DescribeOption{progbarunits}
% \DescribeDefault{5}
% The level bar is devided into boxes or cicles. The number of these items can be set
% via \lstinline`progbarunits` and defaults to five.
% \DescribeMacro{\setprogbarunits}
% It can also be set via \lstinline`\setprogbarunits`\marg{number}.
% 
% \subsubsection{Miscellaneous Options}
% \DescribeOption{draft, final}
% The class also extends the \lstinline`draft` option of the standard article class.
% The otherwise invisible boxes, e.g. for the side bar, will be lightly shaded, 
% so that they are visible.
% The complementary option is \lstinline`final`. 
% Note that if the final option is found, the draft option will have no effect and 
% will be ignored.
% 
% \DescribeOption{biblatex} 
% The \lstinline`biblatex` option will inject some code into a bibliography environemnt, 
% so that the alingment of the hint column will be better. It is turned off by default.
% If the biblatex package is not loaded (this class will not take care of this), 
% it might lead to problems.
% 
% \DescribeOption{signcv} 
% In some countries it is customary to sign the CV.
% The option \lstinline`signcv` can be used to enable this behavior.
% If it is enabled, a box with a signature will be placed in the lower left corner of the 
% first page, below the sidebar contents.
% The box width is also controlled with the sidebar width option.
% An image with the signature is necessary (see below).
%
% \subsubsection{Void Article Options}
% \DescribeOption{twocolumn}
% \watchout
% Since every page has a specific layout passing the \lstinline`twocolumn` option to the 
% inherited article class will lead to problems. Therefore it has been disabled. 
%
% \DescribeOption{titlepage}
%\watchout
% For the same reason as above, creating a title page doesn't make sense and will clash 
% with other settings of this class. 
% Therefore the \lstinline`titlepage` option has also been disabled.
% 
% \section{Changing Fonts and Text Alignment}
% \DescribeMacro{\setheaderfont}
% \DescribeDefault{\mybackslash{}sffamily\mybackslash{}bfseries}
% By default the font used for the header is made bold and serifs are turned off.
% In order to change that layout, the command 
% \lstinline`\setheaderfont`\marg{font specification}
% can be used. This command should be used in the preamble.
%
% \DescribeMacro{\setentryragged}
% \DescribeMacro{\setentryjustified}
% The default layout of an entry in the CV is justified. 
% This can be changed with \lstinline`\setentryragged` to flush the contents to the left
% and leave the right side ragged (\lstinline`\raggedright`). 
% For symmetry the command \lstinline`\setentryjustified` can be used to unset the aabove.
% Both commands will set the internal variable used in the templates at the end of the 
% preamble, so they can only be used there.
%
% \DescribeMacro{\setsigright}
% \DescribeMacro{\setsigleft}
% The commands \lstinline`\setsigright` and \lstinline`\setsigleft` determine
% whether the name below a signature field is right or left aligned.
% This does not affect where the field itself is placed.
% These commands can also only be used in the preamble.
%
% \subsection{Commands to Store and Access Information}
% Personal information can be stored in variables so that they become accessible to the defined templates. 
%
% \DescribeMacro{\author}
% \DescribeMacro{\insertauthor}
% The \lstinline`\author`\marg{author name} command remains unchanged and cannot be used outside of the preamble.
% The contents of that field is copied to the \lstinline`\insertauthor` variable and is accessible 
% throughout the document that way.
%
% \DescribeMacro{\title}
% \DescribeMacro{\inserttitle}
% Similarly to the above, the \lstinline`\title`\marg{title} remains unchanged, with the same limitations applied.
% The contents can be accessed with the \lstinline`\inserttitle` command.
% There is currently no template using this command, but it is recommended to set it to somethin like
% \emph{Curriculum Vitae} for the \PDF{} fields that might use it.
%
% \DescribeMacro{\position}
% \DescribeMacro{\insertposition}
% The \lstinline`\position`\marg{position} macro is defined in a similar way to the above commands.
% The \lstinline`\insertposition` macro can be used to access the content. 
% It is used in the header box and can be set to contain, as the name suggest, the title of the current position.
% A common usage might also be a one line description of one's profession.
%
% \DescribeMacro{\address}
% \DescribeMacro{\insertaddress}
% The address can either be specified by the direct invocation of \lstinline`\address`\marg{address} and
% is accessible via the \lstinline`\insertaddress` macro.
% It is used in the address template (see sectin~\ref{sec:templates-special}), 
% as well as the header of the letter (see section~\ref{sec:letter-env}).
% \DescribeMacro{\street}
% \DescribeMacro{\location}
% Alternatively it the macros \lstinline`\street`\marg{street} and \lstinline`\location`\marg{location} can be used
% to set the values separately. The \lstinline`\insertaddress` will then be contructed from these two values.
% \DescribeMacro{\insertstreet}
% \DescribeMacro{\insertlocation}
% If set, they can be printed with the corresponding \lstinline`\insertstreet` and \lstinline`\insertlocation` macros.
% \watchout
% If the address is set directly the other values remain unchanged. 
% It is therefore also possible to extend the functionality in the follwoing way:
% \iffalse
%<*example>
% \fi
\begin{lstlisting}
\street{123, Example Road}
\location{New Town}
\address{The Old Shack\\\insertstreet\\\insertlocation}
\end{lstlisting}
% \iffalse
%</example>
% \fi
%
% \DescribeMacro{\email}
% \DescribeMacro{\insertemail}
% The command \lstinline`\email`\marg{email} will wrap the input in a hyperlink, which can
% be accessed through the \lstinline`\insertemail` command.
%
% \DescribeMacro{\phone}
% \DescribeMacro{\mobile}
% \DescribeMacro{\insertphone}
% \DescribeMacro{\insertmobile}
% Telephone numbers can be stored with \lstinline`\phone`\marg{number} and \lstinline`\mobile`\marg{number}
% and accessed with the adjacent commands \lstinline`\insertphone` and \lstinline`\insertmobile`.
% Currently the template only provides templates for these two telephone numbers.
%
% \DescribeMacro{\github}
% \DescribeMacro{\insertgithub}
% The macro \lstinline`\github`\marg{username} provides a convenient way to display the GitHub username,
% while also wrapping it into a hyperlink to the associated profile.
%
% \DescribeMacro{\orcid}
% \DescribeMacro{\insertorcid}
% Likewise for an ORCID account, the \lstinline`\orcid`\marg{ID number} will wrap the displayed ID 
% into a hyperlink to the respective profile.
%
% \DescribeMacro{\signaturefile}
% \DescribeMacro{\insertsigfilename}
% In order to sign the CV and the cover letter, a file name must be provided with the macro
% \lstinline`\signaturefile`\marg{filename}. 
% This file name is stored in \lstinline`\insertsigfilename` and is passed to the command 
% \lstinline`includegraphics` in the corresponding templates.
% \watchout
% For the cover letter it is necessary to set this filename, otherwise it will (currently) lead to an error.
%
% \subsection{Templates}
% \label{sec:templates}
% The templates that are provided with this class should make typesetting the contents of the resumee
% easier as they provide predefined spacing and markup.
% 
% \subsubsection{General Templates for the CV}
% \label{sec:templates-general}
%
% \DescribeMacro{\section}
% \DescribeMacro{\subsection}
% The \lstinline`\section`\marg{text} and \lstinline`\subsection`\marg{text} commands are redefined 
% to have no numbers.
% For the former there will be a bar the length of the hint column, while for the latter there is jsut 
% the same amount of space.
% Both commands will use the primary color and the predefined header fonts.
%
% \DescribeMacro{\cvline}
% The \lstinline`\cvline`\oarg{width}\marg{hint}\marg{description} macro divides an entry line in the CV
% into two parts: a hint column with the reset hint, which can also be set with the optional argument,
% and a description box spanning the remaining width of the line.
% It can be used to set any type of centent as a line in the CV.
%
% \DescribeMacro{\cvitemline}
% The \lstinline`\cvitemline`\oarg{width}\marg{icon}\marg{description} is similarly constructed as 
% the above command.
% The general difference is that the hint coulmn is replaced with an icon column.
% If the optional \myarg{length} parameter is not given, it falls back to the definition from the preamle.
% The line itself is defined consisting of two parboxes, spanning the whole width that is available.
% This command is primarily intended for the sidebar, but it could be used somewhere else.
% It is also used for the specialised templates (see section~\ref{sec:templates-special}).
%
% \DescribeMacro{\progressbar}
% The \lstinline`\progressbar`\oarg{width}\marg{level fraction} typesets a progress bar.
% The height of this bar is defined via the class options or in the preamble.
% The width of the bar can be set via the optional parameter \myarg{length}, however,
% in most use cases it will be automatically calculated depending on where it is used.
% The mandatory argument \myarg{level fraction} should be a number between 0 and 1.
% There are some special templates, which use this command (see section~\ref{sec:templates-special}).
%
% \DescribeMacro{\levelbarsquares}
% \DescribeMacro{\levelbarcircles}
% The macros \lstinline`\levelbarsquares`\oarg{width}\marg{level} and 
% \lstinline`\levelbarcircles`\oarg{width} \marg{level} are also progress indicators, 
% which use squares or circles, respectively.
% The number of items is set via the class options or in the preamble.
% The mandatory argument \myarg{level} should be a number within 0 and the set limit.
% The width can be given as the optional argument \myarg{width} and defaults to the width of the hint column.
% There are currently no templates that use these commands.
% 
% \subsubsection{Special Templates for the CV}
% \label{sec:templates-special}
%
%
% \DescribeMacro{\cventry}
% The \lstinline`\cventry` macro is a more specialised template of the \lstinline`\cvline` macro.
% It requires six mandatory arguments, which of course can be empty. 
% It is intended to be used in the following form:\\
% \lstinline`\cventry`\oarg{width}\marg{when}\marg{what}\marg{who}\marg{where}\marg{opt.}\marg{descr.}.\\
% The entry spans the whole available width, the optional \myarg{width} controls the hint column.
% The first argument \myarg{when} contains the hint, which can be a date for example.
% The second argument \myarg{what} will be set in bold font and usually contains the title of the position,
% or a short job description. This is followed by a comma.
% The third argument \myarg{who} is set in italic font and usually contains the name of the employer,
% which is again followed by a comma.
% The fourth argument \myarg{where}  may contain the location of this job and it is set in normal font style.
% If the fifth argument \myarg{opt.} contains something, it is set in the same style as the preceding element, 
% separated by a comma from it. This could be used to specify a grade in the education section for example. 
% If this argument is empty, the comma is skipped. The line will then be terminated by a full stop.
% The sixth argument \myarg{descr.} can also be empty. 
% If it contains text, it will be set on a new line and can be used as a more detailed description of the entry.
% The font size is reduced to small for this paragraph.
%
% \DescribeMacro{\cvskill}
% The \lstinline`\cvskill`\marg{description}\marg{level} is a template using \lstinline`\cvline`, 
% defining a progress bar to \myarg{level} and a \myarg{description}.
% If it is used in the main document part, the progress bar will be set into the hint column, and will 
% take half of its width and is flushed to the right.
% The description is flushed to the left in the main column.
% If it is used in the side bar, the description will start the line, followed by the progress bar.
% The progress bar will then use a third of the whole line, and the description the rest exclluding son separator space.
%
% \DescribeMacro{\cvlanguage}
% The \lstinline`\cvlanguage`\marg{language}\marg{level} is a wrapper to the above command.
% This was introduced, to allow for the possibility of customisation.
% For example, it could be redefined to use a different indicator. \watchout[numeric format]
% \iffalse
%<*example>
% \fi
\begin{lstlisting}
\renewcommand\cvlanguage[2]{%
    \cvline{\levelbarsquares{#2}}{#1}}
\end{lstlisting}
% \iffalse
%</example>
% \fi
%
% \DescribeMacro{\cvlineaddress}
% \DescribeMacro{\cvlineemail}
% \DescribeMacro{\cvlinephone}
% \DescribeMacro{\cvlinemobile}
% The  special versions of the \lstinline`\cvitemline` macro have a preset icon to it.
% The address (\faMapMarker), email (\faEnvelope), phone (\faPhone), and mobile (\faMobile) templates
% are predefined. 
% These macros are primarily intended for the side bar environment.
% They may be used in the main part, too, but they could look somewhat odd.
% Other lines  can be added with the generic command.
% \iffalse
%<*example>
% \fi
\begin{lstlisting}
\cvitemline{\faTerminal}{bash, zsh, fish}
\end{lstlisting}
% \iffalse
%</example>
% \fi
%
% \DescribeMacro{\orcidiconfilename}
% Since there is no \LaTeX{} package available that includes an Orcid icon, 
% it must be provided with an external file. 
% The command \lstinline`\orcidiconfilename` will look for a file with the name \lstinline`orcidicon.pdf`
% and include this as a graphic. 
% If such a file is not found, it will use a palceholder icon made of the letters.
%
% \DescribeMacro{\cvlineorcid}
% \DescribeMacro{\cvlinegithub}
% The orcid and GitHub macros will use their respective icons, similar to the above.
% The contents of the variables defined in the preamble will go into the description column.
%
% \subsection{The CV Environments}
% \label{sec:cv-env}
% \DescribeEnv{polycvfirstpage}
% The class comes with two environments to set the CV contents.
% The first page is special since it also contains a side bar with additional infomation.
% The environment itself sets up the geometry of the page, and the header box including the
% name and position.
% It will allso set up the footer box to close the page.
% It should be used to set a single page only.
%
% \DescribeEnv{polycvsidebar}
% The side bar is part of the first page environment and should be issued inside it.
% It will produce a box with the specified content on the left side of the page. 
% The main templates of this class will have different definintions inside this environment.
%
% The following code block contains an exemplary usage of the two environments.
%
% \iffalse
%<*example>
% \fi
\begin{lstlisting}
\begin{polycvfirstpage}
  \begin{polycvsidebar}
  \includegraphics[width=1.0\linewidth]{example-image-1x1}\\[2ex]
  
  \section{Personal Details}
  \cvitemline{\faStar}{01. 01. 1984,\newline Birthplace, Country}
  \cvlineaddress
  \cvlineemail
  
  \section{Languages}
  \cvlanguage{English (native)}{1}
  
  % etc.
  \end{polycvsidebar}%
  \section{Experience}
  \cventry{years}
    {jobtitle}{employer}{place}
    {optional: ...}
    {optional: job description}

  \section{Skills}
  \cvskill{\LaTeX}{0.8}
  
  % etc.
\end{polycvfirstpage}
\end{lstlisting}
% \iffalse
%</example>
% \fi
% 
% \DescribeEnv{polycvpage}
% The generic CV page environment is similar to the firstpage environment, 
% but it doesn't include space for the ssside bar.
% The margins aare adjusted to fill the whole available width.
% It is useful to typeset a second (or backside) page of the CV.
% The header box of this page will be of the same size as for the first page, but remain empty.
%
% The following codeblock briefly summarises an exaamplle usage:
%
% \iffalse
%<*example>
% \fi
\begin{lstlisting}
\begin{polycvpage}
  \section{Education}
  \cventry{years}
    {degree}{institution}{place}
    {optional: grade/...}
    {optional: comment/ description}

  % etc.
\end{polycvpage}
\end{lstlisting}
% \iffalse
%</example>
% \fi
% 
% \subsection{The Cover Letter Environment}
% \label{sec:letter-env}
% \DescribeEnv{polycvletter}
% The letter environment takes one optional and two mandatory arguments:
% The optional argument is the author's address, which, if unspecified, will
% be constructed from variables set in the preamble.
% The first mandatory argument is the subject line, while the second is
% the recipient's address.
%
% \iffalse
%<*example>
% \fi
\begin{lstlisting}
  \begin{polycvletter}
    [(*@\myarg{Author Address}@*)]{(*@\myarg{Subject line}@*)}{(*@\myarg{Recipient Address}@*)}
    %
    % Letter content
    %
  \end{polycvletter}
\end{lstlisting}
% \iffalse
%</example>
% \fi
%
% At the end of the environment the signature, a file name provided in the preample
% with \lstinline`\signaturefile{<image file>}`, along with the author name will be set.
%
% \DescribeMacro{\opening}
% \DescribeMacro{\closing}
% The letter environment uses indentation specified with the \lstinline`letterindent` 
% option (see section \ref{sec:page-layout}). This looks odd for the opening and closing statement, 
% therefore a wrapper is provided that turns of the indentation. They each take one 
% argument, like  \lstinline`\opening`\marg{statement}.
% These commands will be set in place, i.e. you could include text before and after them.
%
% \DescribeMacro{\ps}
% Occasionally you would like to inlcude a postscript note. 
% The contents of such a note can be stored with \lstinline`\ps`\marg{note} and will be
% produced below the signature of the letter.
% The fontsize is also smaller than the one of the text.
% This command can only be used once; a later occurrence will overwrite a previous one.
%
% \DescribeMacro{\setpsmark} \DescribeDefault{P.S.\~}
% How the postscript is introduced can be changed with \lstinline`\setpsmark`\marg{psmark}.
%
% The following is a minimal working example for a letter created with the polycv class.
%
% \iffalse
%<*example>
% \fi
\begin{lstlisting}
\documentclass[primary=black,14pt,a4paper]{polycv}
\usepackage{mwe} % for this example
\title{Curriculum Vitae}
\author{Example Author}
\signaturefile{example-image-16x9}

\begin{document}
\begin{polycvletter}
  [Address line 1\\ Address line 2]
  {Subject line}
  {Company\\ Address line 1\\ Address line 2}
  \opening{Dear Example,}

  \blindtext[1-3]

  \closing{Sincerely,}
  \ps{\blindtext[1]}
\end{polycvletter}
\end{document}
\end{lstlisting}
% \iffalse
%</example>
% \fi
%
% \subsection{Including external \PDF{} as Appendix}
%
% \DescribeMacro{\pdfappendix} \DescribeDefault{pages=\{1-\}}
% It is occasionally necessary to include certificates, letters of recommendation, 
% or other external sources into an application. 
% This class provides a wrapper to the \lstinline`\includepdf` command from 
% \lstinline`pdfpages`.
% The difference to the original command (which still could be used) is that before 
% including the document, the wrapper unsets the page color. 
% This is necesssaary, even if it is defined as white, due to the way the page color 
% is implemented.
% You can include pages from a \PDF{} document to use it as appendix with the command
% \lstinline`\pdfappendix`\oarg{kwargs}\marg{filename}, where the optional argument accepts
% arguments like \lstinline`pages=1-3` or \lstinline`angle=90`.
% See the \href{https://ctan.org/pkg/pdfpages}{%
%   Documentation of the \lstinline`pdfpages` package}.
%
% \StopEventually{}
%
% \clearpage
% \interfootnotelinepenalty=10000
%
% \section{Implementation}
%
% \subsection{Option definitions}
%
% These simple definitions for key-value pairs as class options are done with kvoptions as
% described in TUGBoat.\footnote{J. Wright, C. Feuersänger, TUGBoat, Vol. 30 (2009), No. 1, p. 110-122.}
% Using the prefix \lstinline`polycv@` as an inbetween to distinguish values set from class options
% in the document.
%
%    \begin{macrocode}
\RequirePackage{kvoptions}
\SetupKeyvalOptions{
  family=polycv,
  prefix=polycv@
}
%    \end{macrocode}
%
% Define options and preset margins for the page setup via the geometry package.
%
%    \begin{macrocode}
\DeclareStringOption[1.0cm]{margins}
\DeclareStringOption[3.0cm]{headerheight}
\DeclareStringOption[0.5cm]{footerheight}
\DeclareStringOption[5.0cm]{sidebarwidth}
\DeclareStringOption[2.7cm]{hintcol}
\DeclareStringOption[0.2cm]{hintcolsep}
\DeclareStringOption[0.8cm]{iconspace}
\DeclareStringOption[1.5ex]{progbarheight}
%    \end{macrocode}
%
% The progressbar, which contains of individual items, can be customised to include a certain
% number of elements. Define the option for this number and set it to five as default.
%
%    \begin{macrocode}
\DeclareStringOption[5]{progbarunits}
%    \end{macrocode}
%
% Define the options for the colors which should be used throughout the document.
% The mixing option is used together with the shading color, to modify the primary color.
%
%    \begin{macrocode}
\DeclareStringOption[blue]{primary}
\DeclareStringOption[white]{secondary}
\DeclareStringOption[orange]{highlight}
\DeclareStringOption[black]{shade}
\DeclareStringOption[10]{mixing}
%    \end{macrocode}
% 
% Define an option for a draft mode, analogous to the standard classes. This will be extended, see below.
% 
%    \begin{macrocode}
\DeclareBoolOption{draft}
\DeclareComplementaryOption{final}{draft}
%    \end{macrocode}
% 
% Define a biblatex mode, which will adjust some space settings in the bibliography, if desired.
%
%    \begin{macrocode}
\DeclareBoolOption{biblatex}
%    \end{macrocode}
%
% Define an option that lets the user pass settings to xcolor, i.e. (the default) \lstinline`dvipsnames`.
%
%    \begin{macrocode}
\DeclareStringOption[dvipsnames]{xcolor}
%    \end{macrocode}
% 
% Define the option whether the CV (front page) is signed or not.
%
%    \begin{macrocode}
\DeclareBoolOption{signcv}
%    \end{macrocode}
%
% Define an option to set the indentation in the letter and let it default to the width or one m.
%
%    \begin{macrocode}
\DeclareStringOption[1em]{letterindent}
%    \end{macrocode}
%
% Warn about options which will not and cannot be used or passed on, because of potential clashes for thiss class.
%
%    \begin{macrocode}
\DeclareVoidOption{twocolumn}{%
  \ClassError{polycv}{%
    Option 'twocolumn' is incompatible because of the special page setup.}}
\DeclareVoidOption{titlepage}{%
  \ClassWarning{polycv}{%
    Option 'titlepage' is incompatible because of the special page setup.}}
%    \end{macrocode}
%
% Pass all other options to the satndart article class.
%
%    \begin{macrocode}
\DeclareDefaultOption{%
  \PassOptionsToClass{\CurrentOptionKey}{article}}
%    \end{macrocode}
%
% Finally, process these options and load the standard class.
%   
%    \begin{macrocode}
\ProcessKeyvalOptions{polycv}
\LoadClass{article}
%    \end{macrocode}
%
% \subsection{Necessary packages}
%
% Some packages are essential for this class and need to be loaded.
%
% \begin{itemize}
%   \item \texttt{kvoptions} is necessary for the use of key-value-options (see above)
%   \item \texttt{etoolbox} is necessary for additional hooks like \lstinline`\AtEndPreamble`.
%   \item \texttt{calc} is necessary to do mathematics with lengths
%   \item \texttt{hyperref} provides hyperlinks within the CV and letter
%   \item \texttt{xcolor} provides convenient color definitions, since hyperref is loaded, this option is set by default.
%     Other options can be added with the class options switch.
%     This package must be loaded before tikz, otherwise there will be an option clash.
%   \item \texttt{tikz} is used for headerbox, symbols, etc.. (must be loaded after xcolor)
%   \item \texttt{graphicx} is used for the inclusion of graphics, e.g. signature and profile picture.
%   \item \texttt{parskip} turns of indentation and improves spacing
%   \item \texttt{geometry} is used for the page layout
%   \item \texttt{fontawesome} provides fancy symbols for the sidebar
%   \item \texttt{pdfpages} are used to include a pdf in the  appendix
% \end{itemize}
%    \begin{macrocode}
\RequirePackage{etoolbox}
\RequirePackage{calc}
\RequirePackage{hyperref}
\PassOptionsToPackage{\polycv@xcolor,hyperref}{xcolor}
\RequirePackage{xcolor}
\RequirePackage{tikz}
\RequirePackage{graphicx}
\RequirePackage[indent=0pt]{parskip}
\RequirePackage{geometry}
\RequirePackage{fontawesome}
\RequirePackage{pdfpages}
%    \end{macrocode}
%
% \subsection{Global settings}
%
% Issue an informative statement that the CV is signed. (To Do.)
%
%    \begin{macrocode}
  \AtEndPreamble{%
    \ifpolycv@signcv%
      \ClassInfo{polycv}{%
        Option 'signcv' enabled; CV will be signed.}%
    \fi
    }%
%    \end{macrocode}
% Set the page style to remove pagenumbers.
%
%    \begin{macrocode}
\pagestyle{empty}
%    \end{macrocode}
%
% \subsubsection{Colors}
%
% Apart from the switches as options, define some commands that can be used to set the colors for the document.
%
% \begin{macro}{\setcolprimary}
% The primary color used for headers.
%    \begin{macrocode}
\newcommand*{\setcolprimary}[1]{%
  \renewcommand*{\polycv@primary}{#1}}
%    \end{macrocode}
% \end{macro}
%
% \begin{macro}{\setcolsecondary}
% The secondary color is the background color, i.e. the page color.
%    \begin{macrocode}
\newcommand*{\setcolsecondary}[1]%
  {\renewcommand*{\polycv@secondary}{#1}}
%    \end{macrocode}
% \end{macro}
%
% \begin{macro}{\setcolhighlight}
% If text should be highlighted another color can be set.
%    \begin{macrocode}
\newcommand*{\setcolhighlight}[1]%
  {\renewcommand*{\polycv@highlight}{#1}}
%    \end{macrocode}
% \end{macro}
%
% \begin{macro}{\setcolshade}
% The color which is used for shading (damping) the primary color.
%    \begin{macrocode}
\newcommand*{\setcolshade}[1]%
  {\renewcommand*{\polycv@shade}{#1}}
%    \end{macrocode}
% \end{macro}
%
% Disable these changing color commands for body, then define the colors in internal
% macros along with some mixing.
% Then set the text color and the color for the page for the entire document.
%
%    \begin{macrocode}
\@onlypreamble\setcolprimary
\@onlypreamble\setcolsecondary
\@onlypreamble\setcolhighlight
\@onlypreamble\setcolshade
\AtEndPreamble{%
  \colorlet{polycvpagecolor}{\polycv@secondary}
  \colorlet{polycvheaderbackground}{\polycv@primary}
  \colorlet{polycvheadertext}{\polycv@secondary}
  \colorlet{polycvsectioncolor}{\polycv@primary}
  \colorlet{polycvtextcolor}{%
    \polycv@primary!\polycv@mixing!\polycv@shade}
  \colorlet{polycvhighlight}{\polycv@highlight}
  \colorlet{polycvsidebarbackground}{\polycv@secondary}
  \colorlet{polycvprogbarfilled}{\polycv@primary}
  \colorlet{polycvprogbarnofill}{%
    \polycv@primary!10!\polycv@secondary}
}
\AfterPreamble{%
  \color{polycvtextcolor}\pagecolor{polycvpagecolor}}
%    \end{macrocode}
%
% \subsubsection{Extending draft mode}
% This section emulates and extends the draft mode.
% It writes an information message to the log.
% It then sets the overfull boxes identifier.\footnote{%
% \textit{How to test whether report has been called with “draft” option}%
% \url{https://tex.stackexchange.com/a/240128/33413}}
% Passes the \lstinline`showframe` option to the geometry package to make 
% the page borders visible.
% It also sets the color of the side bar, to make the tikz boxes visible in draft mode.
%    \begin{macrocode}
\ifpolycv@draft%
  \ClassInfo{polycv}{Option 'draft' is enabled}
  \setlength\overfullrule{1em}
  \PassOptionsToPackage{showframe}{geometry}
  \AtEndPreamble{%
    \colorlet{polycvsidebarbackground}{%
      \polycv@primary!20!\polycv@secondary}%
  }
\fi
%    \end{macrocode}
% \subsubsection{Biblatex Mode}
% In this option the biblatex package is set to use the same width as the date column 
% for the numeration.
%    \begin{macrocode}
\ifpolycv@biblatex%
  \ClassInfo{polycv}{Option 'biblatex' is enabled}
  \AtEndPreamble{%
    \setlength{\biblabelsep}{\polycvhintcolsep}%
    \DeclareFieldFormat{labelnumberwidth}{%
      \makebox[\polycvhintcol][r]{[#1]}}%
  }%
\fi
%    \end{macrocode}
%
% \subsubsection{Lengths}
%
% Define the lengths and use defaults or parameters from the class options for the layout.
%
%    \begin{macrocode}
\newlength{\polycvheaderheight}
\setlength{\polycvheaderheight}{\polycv@headerheight}
\newlength{\polycvfooterheight}
\setlength{\polycvfooterheight}{\polycv@footerheight}
\newlength{\polycvmargins}
\setlength{\polycvmargins}{\polycv@margins}
\newlength{\polycvsidebarwidth}
\setlength{\polycvsidebarwidth}{\polycv@sidebarwidth}
\newlength{\polycvhintcol}
\setlength{\polycvhintcol}{\polycv@hintcol}
\newlength{\polycvhintcolsep}
\setlength{\polycvhintcolsep}{\polycv@hintcolsep}
\newlength{\polycviconspace}
\setlength{\polycviconspace}{\polycv@iconspace}
\newlength{\polycvprogbarheight}
\setlength{\polycvprogbarheight}{\polycv@progbarheight}
\newlength{\polycvletterindent}
\setlength{\polycvletterindent}{\polycv@letterindent}
%    \end{macrocode}
%
% The lengths can also be altered with the following commands.
% \begin{macro}{\setheaderheight}
% Define the hight for the header box.
%    \begin{macrocode}
\newcommand*{\setheaderheight}[1]{%
  \setlength{\polycvheaderheight}{#1}}
%    \end{macrocode}
% \end{macro}
%
% \begin{macro}{\setfooterheight}
% Define the height for the footer box.
%    \begin{macrocode}
\newcommand*{\setfooterheight}[1]{%
  \setlength{\polycvfooterheight}{#1}}
%    \end{macrocode}
% \end{macro}
%
% \begin{macro}{\setmargins}
% Define the width of the margins.
%    \begin{macrocode}
\newcommand*{\setmargins}[1]{%
  \setlength{\polycvmargins}{#1}}
%    \end{macrocode}
% \end{macro}
%
% \begin{macro}{\sethintcolwidth}
% Define the hint column width (where dates and level, etc. are set).
%    \begin{macrocode}
\newcommand*{\sethintcolwidth}[1]{%
  \setlength{\polycvhintcol}{#1}}
%    \end{macrocode}
% \end{macro}
%
% \begin{macro}{\sethintcolsep}
% Define the space between the hint column and the description column.
%    \begin{macrocode}
\newcommand*{\sethintcolsep}[1]{%
  \setlength{\polycvhintcolsep}{#1}}
%    \end{macrocode}
% \end{macro}
%
% \begin{macro}{\seticonspace}
% Define the space reserved for the items in the sidebar box.
%    \begin{macrocode}
\newcommand*{\seticonspace}[1]{%
  \setlength{\polycviconspace}{#1}}
%    \end{macrocode}
% \end{macro}
%
% \begin{macro}{\setprogbarheight}
% Define the standard height of the progress bar. This will also apply to the levels bar.
%    \begin{macrocode}
\newcommand*{\setprogbarheight}[1]{%
  \setlength{\polycvprogbarheight}{#1}}
%    \end{macrocode}
% \end{macro}
%
% These commands need to be disabled for the body, 
% so that automatically derived lengths will still be useable.
%    \begin{macrocode}
\@onlypreamble\setheaderheight
\@onlypreamble\setfooterheight
\@onlypreamble\setmargins
\@onlypreamble\sethintcolwidth
\@onlypreamble\sethintcolsep
\@onlypreamble\seticonspace
\@onlypreamble\setprogbarheight
%    \end{macrocode}
% 
% Additional lengths, which are derived from the above,
% will be set at the end of the preamble.
%
%    \begin{macrocode}
\newlength{\polycvleftmargin}
\newlength{\polycvtopmargin}
\newlength{\polycvbottommargin}
\AtEndPreamble{%
  \setlength{\polycvleftmargin}{%
    \polycvsidebarwidth+2\polycvmargins}
  \setlength{\polycvtopmargin}{%
    \polycvheaderheight+\polycvmargins}
  \setlength{\polycvbottommargin}{%
    \polycvfooterheight+\polycvmargins}
}
%    \end{macrocode}
%
% \subsubsection{Fonts and Styles}
%
% \begin{macro}{\polycv@headerfont}
% \begin{macro}{\setheaderfont}
% Define the font family for the header and make it bold (by default) and store
% it in a variable. Provide access to this cariable via a set command.
%    \begin{macrocode}
\providecommand\polycv@headerfont{}
\renewcommand{\polycv@headerfont}{\sffamily\bfseries}
\providecommand\setheaderfont{}
\renewcommand*{\setheaderfont}[1]{%
  \renewcommand*{\polycv@headerfont}{#1}}
\@onlypreamble\setheaderfont
%    \end{macrocode}
% \end{macro}
% \end{macro}
%
% \begin{macro}{\polycv@entry@ragged}
% \begin{macro}{\setentryragged}
% \begin{macro}{\setentryjustified}
% Provide a variable to store how the output is set, i.e. the standard raggedness.
% Provide commands to change the raggedness of the output. 
% These need to be set in the preamble, so they are disabled for the body.
%    \begin{macrocode}
\providecommand\polycv@entry@ragged{}
\newcommand*{\setentryragged}{%
  \renewcommand{\polycv@entry@ragged}{\raggedright}}
\newcommand*{\setentryjustified}{%
  \renewcommand{\polycv@entry@ragged}{}}
\@onlypreamble\setentryragged
\@onlypreamble\setentryjustified
%    \end{macrocode}
% \end{macro}
% \end{macro}
% \end{macro}
%
% \begin{macro}{\polycv@sig@align}
% \begin{macro}{\setsigright}
% \begin{macro}{\setsigleft}
% Define the alignment for the signature, which is stored in a variable.
%  Provide commands to change it, and disable them for the body.
%    \begin{macrocode}
\providecommand\polycv@sig@align{right}
\newcommand*{\setsigright}{%
  \renewcommand{\polycv@sig@align}{right}}
\newcommand*{\setsigleft}{%
  \renewcommand{\polycv@sig@align}{left}}
\@onlypreamble\setsigright
\@onlypreamble\setsigleft
%    \end{macrocode}
% \end{macro}
% \end{macro}
% \end{macro}
%
% \subsection{Macros and Environments}
%
% Define special commands and environment.s
%
% \begin{macro}{\setprogbarunits}
% Control the number of items for the levels bar.
%    \begin{macrocode}
\newcommand*{\setprogbarunits}[1]{%
  \renewcommand*{\polycv@progbarunits}{#1}}
\@onlypreamble\setprogbarunits
%    \end{macrocode}
% \end{macro}
%
% \subsubsection{Personal Information}
%
% This section defines the variables, which will store the personal information.
% These are used in the template lines of the sidebar.
%
% \begin{macro}{\insertauthor}
% \begin{macro}{\inserttitle}
% The author and title commands already exist, so they only need to be stored.
%    \begin{macrocode}
\providecommand*\insertauthor{}
\renewcommand{\insertauthor}{\@author}
\providecommand\inserttitle{}
\renewcommand{\inserttitle}{\@title}
%    \end{macrocode}
% \end{macro}
% \end{macro}
%
% \begin{macro}{\insertposition}
% \begin{macro}{\position}
% Store the position, which is used as the sub-heading in the title banner.
% Disable this command because it is needed before the header is set.
%    \begin{macrocode}
\providecommand*\insertposition{}
\providecommand*\position{}
\renewcommand{\position}[1]{%
  \renewcommand{\insertposition}{#1}}
\@onlypreamble\position
%    \end{macrocode}
% \end{macro}
% \end{macro}
%
% \begin{macro}{\cvitemline}
% This defines a generic line consisting of two parboxes.
% This command is primarily needed for the sidebar, but it could be used somewhere else.
% The left parbox is used as a hint column and contains an icon,
% the right the value of a variable.
%    \begin{macrocode}
\providecommand\cvitemline[3][\polycviconspace]{%
  \parbox[t]{#1}{\centering#2}%
  \parbox[t]{\linewidth-#1}{#3}\par\vspace{0.3\baselineskip}}
%    \end{macrocode}
% \end{macro}
%
% \begin{macro}{\insertaddress}
% \begin{macro}{\insertstreet}
% \begin{macro}{\insertlocation}
% Variables to store an address.  
%    \begin{macrocode}
\providecommand\insertaddress{}
\providecommand\insertstreet{}
\providecommand\insertlocation{}
%    \end{macrocode}
% \end{macro}
% \end{macro}
% \end{macro}
%
% \begin{macro}{\address}
% \begin{macro}{\street}
% \begin{macro}{\location}
% Commands to set the values used for the address
%    \begin{macrocode}
\providecommand\address{}
\providecommand\street{}
\providecommand\location{}
\renewcommand{\address}[1]{%
  \renewcommand{\insertaddress}{#1}}
\renewcommand{\street}[1]{%
  \renewcommand{\insertstreet}{#1}}
\renewcommand{\location}[1]{%
  \renewcommand{\insertlocation}{#1}}
%    \end{macrocode}
% Build the address for street and location.
%    \begin{macrocode}
\renewcommand{\insertaddress}{\insertstreet\\\insertlocation}
%    \end{macrocode}
% \end{macro}
% \end{macro}
% \end{macro}
%
% \begin{macro}{\cvlineaddress}
% Template to produce the address for the sidebar.
%    \begin{macrocode}
\providecommand{\cvlineaddress}{%
  \cvitemline{\faMapMarker}{\insertaddress}}
%    \end{macrocode}
% \end{macro}
%
% \begin{macro}{\insertemail}
% \begin{macro}{\email}
% \begin{macro}{\cvlineemail}
% Store and set email, template for the sidebar.
%    \begin{macrocode}
\providecommand*\insertemail{}
\providecommand*\email{}
\renewcommand{\email}[1]{%
  \renewcommand{\insertemail}{\href{mailto:#1}{#1}}}
\providecommand{\cvlineemail}{%
  \cvitemline{\faEnvelope}{\insertemail}}
%    \end{macrocode}
% \end{macro}
% \end{macro}
% \end{macro}
%
% \begin{macro}{\insertphone}
% \begin{macro}{\phone}
% \begin{macro}{\cvlinephone}
% Store and set phone number, template for the sidebar.
%    \begin{macrocode}
\providecommand*\insertphone{}
\providecommand*\phone{}
\renewcommand{\phone}[1]{\renewcommand{\insertphone}{#1}}
\providecommand{\cvlinephone}{%
  \cvitemline{\faPhone}{\insertphone}}
%    \end{macrocode}
% \end{macro}
% \end{macro}
% \end{macro}
%
% \begin{macro}{\insertmobile}
% \begin{macro}{\mobile}
% \begin{macro}{\cvlinemobile}
% Store and set mobile number, template for the sidebar.
%    \begin{macrocode}
\providecommand*\insertmobile{}
\providecommand*\mobile{}
\renewcommand{\mobile}[1]{\renewcommand{\insertmobile}{#1}}
\providecommand{\cvlinemobile}{%
  \cvitemline{\faMobile}{\insertmobile}}
%    \end{macrocode}
% \end{macro}
% \end{macro}
% \end{macro}
%
% \begin{macro}{\insertgithub}
% \begin{macro}{\github}
% \begin{macro}{\cvlinegithub}
% Store and set username for GitHub, template for the sidebar.
%    \begin{macrocode}
\providecommand*\insertgithub{}
\providecommand*\github{}
\renewcommand{\github}[1]{\renewcommand{\insertgithub}{%
  \href{https://github.com/#1}{github.com/#1}}}
\providecommand{\cvlinegithub}{%
  \cvitemline{\faGithub}{\insertgithub}}
%    \end{macrocode}
% \end{macro}
% \end{macro}
% \end{macro}
%
% \begin{macro}{\orcidiconfilename}
% This defines the orcid icon either from the file \lstinline`orcidicon.pdf` 
% if it is found by \LaTeX.
% I am unaware of a \LaTeX package that defines this symbol.\footnote{%
% For more details see:  
% \textit{Is there a standard way to include ORCID in TeX / PDF?}
% (\url{https://tex.stackexchange.com/q/275578/33413})}
% If such a file is  not found, it will place as small square of bold letters spelling orcid.
%    \begin{macrocode}
\providecommand*\orcidiconfilename{}
\IfFileExists{orcidicon.pdf}{%
  \renewcommand{\orcidiconfilename}{%
    \raisebox{-0.2em}{%
      \includegraphics[height=1.0em]{orcidicon.pdf}}}%
}{%
  \renewcommand{\orcidiconfilename}{%
    \raisebox{-0.2em}{%
      \resizebox{1em}{1em}{\textbf{ORCID}}}}%
}
%    \end{macrocode}
% \end{macro}
%
% \begin{macro}{\insertorcid}
% \begin{macro}{\orcid}
% \begin{macro}{\cvlineorcid}
% Store and set the orcid number, template for the sidebar.
%    \begin{macrocode}
\providecommand*\insertorcid{}
\providecommand*\orcid{}
\renewcommand{\orcid}[1]{\renewcommand{\insertorcid}{%
  \href{http://orcid.org/#1}{orcid.org/#1}}}
\providecommand{\cvlineorcid}{%
  \cvitemline{\orcidiconfilename}{\insertorcid}}
%    \end{macrocode}
% \end{macro}
% \end{macro}
% \end{macro}
%
% \begin{macro}{\insertsigfilename}
% \begin{macro}{\signaturefile}
% Store the filename of an image used to produce the signature.
%    \begin{macrocode}
\providecommand\insertsigfilename{}
\providecommand\signaturefile{}
\renewcommand{\signaturefile}[1]{%
  \renewcommand{\insertsigfilename}{#1}}
%    \end{macrocode}
% \end{macro}
% \end{macro}
%
% \subsubsection{Redefinition of Section Headers}
%
% \begin{macro}{\section}
% \begin{macro}{\subsection}
% The section and subsection commands are redefined to have no numbers.
% Instead there will be a bar the length of the hint column, and they will use the primary color.
%
%    \begin{macrocode}
\renewcommand{\section}[1]{%
  \ignorespaces%
  \parbox[b]{1\linewidth}{%
    \color{polycvsectioncolor}%
    \Large\polycv@headerfont%
    \rule{\polycvhintcol}{1ex}\hspace{\polycvhintcolsep}%
    \parbox{1.0\linewidth-\polycvhintcol-\polycvhintcolsep}{%
      #1%
    }% 
  }\vspace{0.5\baselineskip}}
\renewcommand{\subsection}[1]{%
  \parbox[b]{1\linewidth}{%
    \color{polycvsectioncolor}%
    \polycv@headerfont%
    \hspace{\polycvhintcol}\hspace{\polycvhintcolsep}%
    \parbox{1.0\linewidth-\polycvhintcol-\polycvhintcolsep}{%
      #1%
    }% 
  }\vspace{0.3\baselineskip}}
%    \end{macrocode}
% \end{macro}
% \end{macro}
%
% \subsubsection{Header and Footer Boxes}
%
% \begin{macro}{\polycv@header}
% The header box will be used internally only. 
% It is a simple box produced with tikz, and uses the author and position variables.
% It is typeset in the primary color of the template.
%
%    \begin{macrocode}
\newcommand{\polycv@header}[3][\polycvheaderheight]{%
  \begin{tikzpicture}[remember picture,overlay]%
    \node [rectangle, %
           fill           = polycvheaderbackground, %
           anchor         = north, %
           minimum width  = \paperwidth, %
           minimum height = #1%
          ] (headerbox) at (current page.north){};%
    \node [anchor = mid] (cv-name) at (headerbox) {%
      \Huge\polycv@headerfont\color{polycvheadertext}\textbf{%
        #2
      }%
    };%
    \node [anchor = north] at (cv-name.south) {%
      \Large\polycv@headerfont\color{polycvheadertext}%
      #3%
    };%
  \end{tikzpicture}%
  \ignorespaces%
}
%    \end{macrocode}
% \end{macro}
%
% \begin{macro}{\polycv@footer}
% The footer box will be used internally only. 
% It is a simple box produced with tikz, which can be used to ddisplay page numbers.
% If the argument to the macro is empty, nothing will be printed, 
% otherwise dashes are used to surround the command.
% It is typeset in the primary color of the template.
%
%    \begin{macrocode}
\newcommand{\polycv@footer}[2][\polycvfooterheight]{%
  \begin{tikzpicture}[remember picture,overlay]%
    \node [rectangle, %
           fill           = polycvheaderbackground, %
           anchor         = south, %
           minimum width  = \paperwidth, %
           minimum height = #1%
          ] (footerbox) at (current page.south){};%
    \node [anchor = center] (pagenumber) at (footerbox) {%
      \polycv@headerfont\color{polycvheadertext}%
      \if\relax\detokenize{#2}\relax%
      \else%
      - #2 -%
      \fi%
    };%
  \end{tikzpicture}%
  \ignorespaces%
}
%    \end{macrocode}
% \end{macro}
%
% \begin{macro}{\polycv@signature}
% The signature box will be used internally only. 
% It is a simple box produced with tikz, which wraps around an image defined earlier.
% It will also print the author name below the image.
% By default the box has the same background color as the page.
%
%    \begin{macrocode}
\newcommand{\polycv@signature}[2][\polycvsidebarwidth]{%
  \begin{tikzpicture}[remember picture, overlay]%
    \node [rectangle, %
           anchor = south west, %
           text width = #1, %
           align = \polycv@sig@align, %
           minimum height = 0.75#1, %
           yshift = \polycvbottommargin, %
           xshift = \polycvmargins, %
           fill = polycvsidebarbackground % 
          ] (cv-signbox) at (current page.south west) {%
      \includegraphics[width=#1]{#2}\\%
      \insertauthor%
    };
  \end{tikzpicture}%
  \ignorespaces%
}
%    \end{macrocode}
% \end{macro}
%
% \subsubsection{Level and Progress Bar}
%
% First a few simple filled and unfilled symbols are defined as tikz pictures.
% \begin{macro}{\polycv@square@nofill}
% \begin{macro}{\polycv@square@filled}
% Square symbols, the height and width are set to the height of the progressbar, 
% which is defined via the options.
%
%    \begin{macrocode}
\newcommand{\polycv@square@nofill}[1][\polycvprogbarheight]{%
  \begin{tikzpicture}[x=#1, y=#1]%
    \filldraw[polycvprogbarnofill] (0,0) rectangle (1,1);% 
    \draw[polycvtextcolor, thick] (0,0) rectangle (1,1);% 
  \end{tikzpicture}}
\newcommand{\polycv@square@filled}[1][\polycvprogbarheight]{%
  \begin{tikzpicture}[x=#1, y=#1]%
    \filldraw[polycvprogbarfilled] (0,0) rectangle (1,1);% 
    \draw[polycvtextcolor, thick] (0,0) rectangle (1,1);% 
  \end{tikzpicture}}
%    \end{macrocode}
% \end{macro}
% \end{macro}
%
% \begin{macro}{\polycv@circle@nofill}
% \begin{macro}{\polycv@circle@filled}
% Analogous circle symbols.
%    \begin{macrocode}
\newcommand{\polycv@circle@nofill}[1][\polycvprogbarheight]{%
  \begin{tikzpicture}[x=#1, y=#1]%
    \filldraw[polycvprogbarnofill] (0.5,0.5) circle [radius=0.5];% 
    \draw[polycvtextcolor, thick] (0.5,0.5) circle [radius=0.5];% 
  \end{tikzpicture}}
\newcommand{\polycv@circle@filled}[1][\polycvprogbarheight]{%
  \begin{tikzpicture}[x=#1, y=#1]%
    \filldraw[polycvprogbarfilled] (0.5,0.5) circle [radius=0.5];% 
    \draw[polycvtextcolor, thick] (0.5,0.5) circle [radius=0.5];% 
  \end{tikzpicture}}
%    \end{macrocode}
% \end{macro}
% \end{macro}
%
% \begin{macro}{\polycv@repeating}
% This is a helper function to repeat typesetting a symbol. 
% It will fill the spaces between the symbols.
%
%    \begin{macrocode}
\newcommand{\polycv@repeating}[2][5]{%
  \ifnum#1 > 0%
    \newcount\rep%
    \rep0%
    \loop\ifnum\rep < \number\numexpr#1-1\relax%
      #2\hfill%
      \advance\rep by 1% 
    \repeat%
    #2%
  \fi%
}
%    \end{macrocode}
% \end{macro}
%
% \begin{macro}{\polycv@level@bar}
% The level bar is also a helper function, which wraps around the repeating function.
% It sets its contents as a parbox with the width of the hint column (\#1 by default).
% It also needs a level to display (\#2), which is the number of filled elements (\#3).
% It the computes from the options of total units the number of unfilled elements(\#4)
% to typeset.
%
%    \begin{macrocode}
\newcommand{\polycv@level@bar}[4][\polycvhintcol]{%
  \newcount\total@units%
  \total@units\polycv@progbarunits%
  \newcount\level%
  \ifnum#2 < \total@units \level#2 \else \level\total@units \fi%
  \newcount\elevel%
  \elevel\numexpr\total@units-\level\relax%
  \def\symbol@filled{#3}%
  \def\symbol@nofill{#4}%
  \parbox{#1}{%
    \ifnum\level = 0%
      \polycv@repeating[\total@units]{\symbol@nofill}%
    \else 
      \ifnum\level = \total@units%
        \polycv@repeating[\total@units]{\symbol@filled}%
      \else
        \polycv@repeating[\level]{\symbol@filled}%
        \hfill%
        \polycv@repeating[\elevel]{\symbol@nofill}%
      \fi%
    \fi%
  }%
}
%    \end{macrocode}
% \end{macro}
%
% \begin{macro}{\levelbarsquares}
% \begin{macro}{\levelbarcircles}
% Wrapper commands as user interfaces to the level bar.
%
%    \begin{macrocode}
\newcommand{\levelbarsquares}[2][\polycvhintcol]{%
  \polycv@level@bar[#1]{#2}{%
    \polycv@square@filled}{\polycv@square@nofill}%
}
\newcommand{\levelbarcircles}[2][\polycvhintcol]{%
  \polycv@level@bar[#1]{#2}{%
    \polycv@circle@filled}{\polycv@circle@nofill}%
}
%    \end{macrocode}
% \end{macro}
% \end{macro}
%
% \begin{macro}{\progressbar}
% The progressbar is a simple tikz image consiting of three rectangles:
% one denoting the level using the fill color, one using the empty color,
% and an empty box encompasing all as a frame.
%  
%    \begin{macrocode}
\newcommand{\progressbar}[2][\polycvhintcol]{%
  \hfill%
  \begin{tikzpicture}[x=0.99*#1, y=\polycvprogbarheight]%
    \filldraw[polycvprogbarfilled] (0, 0) rectangle (#2, 1) ;%
    \filldraw[polycvprogbarnofill] (#2, 0) rectangle (1, 1);%
    \draw[polycvtextcolor, thick] (0, 0) rectangle (1, 1) ;%
  \end{tikzpicture}%
}%
%    \end{macrocode}
% \end{macro}
%
% \subsubsection{Commands for the main Body}
%
% \begin{macro}{\cvline}
%  A generic command to typeset a line in the CV.
%  It consists of two parboxes, the left being the hint column, and the right the description.
%  One example usage is: \lstinline`\cvline[width]{hint}{description}`.
%    \begin{macrocode}
\providecommand\cvline{}
\renewcommand{\cvline}[3][\polycvhintcol]{%
  \par%
  \parbox[t]{#1}{\strut\raggedleft #2}%
  \hspace{\polycvhintcolsep}%
  \parbox[t]{\linewidth-\polycvhintcolsep-#1}{%
    \strut\polycv@entry@ragged #3}%
  \par\vspace{0.2\baselineskip}%
}
%    \end{macrocode}
% \end{macro}
%
% \begin{macro}{\cventry}
% The entry defines a block. It uses the command above, but fills in more details 
% in the description parbox of the template.
% An example usage is: \lstinline`\cventry[width]{when}{what}{who}{where}{optional}{optional}`.
%    \begin{macrocode}
\providecommand\cventry{}
\renewcommand{\cventry}[7][\polycvhintcol]{%
  \cvline[#1]{#2}{%
    \textbf{#3}, \textit{#4}, #5%
    \if\relax\detokenize{#6}\relax%
    \else%
    , #6%
    \fi%
    \if\relax\detokenize{#7}\relax%
    .%
    \else%
    .\newline\small #7%
    \fi%
  }%
}
%    \end{macrocode}
% \end{macro}
%
% \begin{macro}{\cvskill}
% This command uses a progressbar in the hint column.
%    \begin{macrocode}
\providecommand\cvskill[2]{%
  \cvline[\polycvhintcol]{%
    \progressbar[0.5\linewidth]{#2}}{#1}}
%    \end{macrocode}
% \end{macro}
%
% \begin{macro}{\cvlanguage}
% This is basically a wrapper to the above command.
%    \begin{macrocode}
\providecommand\cvlanguage[2]{%
  \cvskill{#1}{#2}}
%    \end{macrocode}
% \end{macro}
%
% \begin{environment}{polycvfirstpage}
% This environment redefines the page layout for the first page, 
% and calls the header and footer templates.
%
%    \begin{macrocode}
\newenvironment{polycvfirstpage}{%
  \newgeometry{%
    left=\polycvleftmargin,
    top=\polycvtopmargin,
    right=\polycvmargins,
    bottom=\polycvbottommargin,
    nohead,nofoot}%
  \polycv@header{\insertauthor}{\insertposition}%
  \polycv@footer{}%
}{%
  \restoregeometry%
  \ignorespaces%
}
%    \end{macrocode}
% \end{environment}
%
% \begin{environment}{polycvpage}
% This environment defines the page layout for a 'regular' page.
% It may cause problems, when the content is overflowing, 
% i.e. producing more than one page.
% Unfortunately, page handling has (still) to be done manually.
% It also calls the header and footer templates, but they are empty by default.
%    \begin{macrocode}
\newenvironment{polycvpage}{%
  \newgeometry{%
    top=\polycvtopmargin,%
    left=2\polycvmargins,%
    right=2\polycvmargins,%
    bottom=\polycvbottommargin,%
    nohead,nofoot}%
  \polycv@header{}{}%
  \polycv@footer{}%
}{%
  \restoregeometry%
  \ignorespaces%
}
%    \end{macrocode}
% \end{environment}
%
% \subsubsection{Side Bar Definitions}
%
% \begin{macro}{\polycv@sidebar@skill}
% This template is used in the sidebar.
% It consists of one parbox spanning the full width of the line,
% encapsulating two more parboxes to set the progressbar and the description.
%
%    \begin{macrocode}
\providecommand\polycv@sidebar@skill[3][\polycvhintcolsep]{%
  \parbox{1.0\linewidth}{%
    \hspace{#1}%
    \parbox[b]{2\linewidth/3-4#1}{#2}%
    \hspace{#1}\hfill%
    \parbox[b]{\linewidth/3}{%
      \progressbar[\linewidth]{#3}%
      }%
    \vspace{0.2\baselineskip}}%
}
%    \end{macrocode}
% \end{macro}
%
% \begin{environment}{polycvsidebar}
% This environment provides a box as a tikz picture for the sidebar.
% It first saves the original definition of the section command, then redefines it for a better fit.
% It also redefines the skill and language commands for a better fit within the box.
% At the end, it restores the original definitions.
% If the option to sign the CV is active, another box will be inserted at the bottom off the paaage.
%
%    \begin{macrocode}
\newenvironment{polycvsidebar}{%
  \let\section\save@section%
  \newcommand{\section}[1]{%
  \parbox[b]{1\linewidth}{%
    \color{polycvsectioncolor}\large\polycv@headerfont{##1}% 
  }\par\vspace{0.5\baselineskip}}%
  \let\cvlanguage\save@polycvlanguage%
  \newcommand{\cvlanguage}[2]{%
    \polycv@sidebar@skill{##1}{##2}}%
  \let\cvskill\save@polycvskill%
  \newcommand{\cvskill}[2]{%
    \polycv@sidebar@skill{##1}{##2}}%
  \begin{tikzpicture}[remember picture, overlay]%
    \node [anchor = north west, %
           text width = \polycvsidebarwidth, %
           xshift = \polycvmargins, %
           yshift = -\polycvheaderheight-\polycvmargins, %
           fill = polycvsidebarbackground %
          ] (polycvsidebarbox) at (current page.north west) %
    \bgroup%
}{%
    \egroup;%
  \end{tikzpicture}%
  \let\save@section\section%
  \let\save@polycvlanguage\cvlanguage%
  \let\save@polycvlanguage\cvskill%
  \ifpolycv@signcv%
    \polycv@signature{\insertsigfilename}%
  \fi%
  \ignorespaces%
}%
%    \end{macrocode}
% \end{environment}
%
% \subsection{Cover Letter Definitions}
%
% In this section the commands for the letter template are set up.
%
% \begin{macro}{\polycv@letter@from}
% The header box is now substituted to include the address. 
% The colourscheme is also adjusted to make the box lesss bulky.
% Therefore the usual background colour and the text color will be used.
%    \begin{macrocode}
\newcommand{\polycv@letter@from}[3][\polycvheaderheight]{%
  \begin{tikzpicture}[remember picture, overlay]%
    \node [rectangle, %
           fill           = polycvheadertext, %
           anchor         = north, %
           minimum width  = \paperwidth, %
           minimum height = #1-1ex%
          ] (headerbox) at (current page.north){};
    \node [rectangle, %
           fill           = polycvheaderbackground, %
           anchor         = north, %
           minimum width  = \paperwidth, %
           minimum height = 1ex%
          ] (headerline) at (headerbox.south){};
    \node [anchor = east, %
           text width = 5cm, %
           xshift     = -2\polycvmargins
          ] (polycv-letter-address) at (headerbox.east) {%
      \polycv@headerfont\color{polycvheaderbackground}%
      \parbox[t]{1.0\linewidth}{%
        \raggedleft%
        \textbf{\Large #2}\\%
        \textit{\normalsize #3}%
        }%
    };%
  \end{tikzpicture}%
  \ignorespaces%
}
%    \end{macrocode}
% \end{macro}
%
% \begin{macro}{\polycv@letter@to}
% This defines the field of the addressee.
%
%    \begin{macrocode}
\newcommand{\polycv@letter@to}[1]{%
  \noindent\parbox[t][4cm][c]{0.4\linewidth}{\raggedright#1}\par%
  \vspace{1\baselineskip}%
}
%    \end{macrocode}
% \end{macro}
%
% \begin{macro}{\polycv@letter@date}
% This defines and formates the date field.
%
%    \begin{macrocode}
\newcommand{\polycv@letter@date}{%
  \noindent\parbox[t]{1.00\linewidth}{\raggedleft\today}\par%
}
%    \end{macrocode}
% \end{macro}
%
% \begin{macro}{\polycv@letter@subject}
% This defines and formates the subject field.
%
%    \begin{macrocode}
\newcommand{\polycv@letter@subject}[1]{%
  \noindent\parbox[t]{1.00\linewidth}{%
    \raggedright\textbf{#1}}%
}
%    \end{macrocode}
% \end{macro}
%
% \begin{macro}{\opening}
% \begin{macro}{\closing}
% This defines and formates the opening and closing remarks
%
%    \begin{macrocode}
\providecommand\opening{}
\renewcommand{\opening}[1]{\noindent{}#1\par}
\providecommand\closing{}
\renewcommand{\closing}[1]{\noindent{}#1\par}
%    \end{macrocode}
% \end{macro}
% \end{macro}
%
% \begin{macro}{\polycv@letter@psmark}
% Define the default postscript mark.
%
%    \begin{macrocode}
\newcommand{\polycv@letter@psmark}{P.S.~}
%    \end{macrocode}
% \end{macro}
%
% \begin{macro}{\setpsmark}
% Command to change the default postscript mark.
%    \begin{macrocode}
\providecommand\setpsmark{}
\renewcommand*{\setpsmark}[1]{%
  \renewcommand{\polycv@letter@psmark}{#1}}
%    \end{macrocode}
% \end{macro}
%
% \begin{macro}{\polycv@letter@ps}
% \begin{macro}{\ps}
% Variable to store the content of the postscript remark and command to set it.
%    \begin{macrocode}
\providecommand\polycv@letter@ps{}%
\providecommand\ps{}
\renewcommand{\ps}[1]{%
  \renewcommand{\polycv@letter@ps}{%
    \noindent\footnotesize\polycv@letter@psmark#1\par%
  }%
}%
%    \end{macrocode}
% \end{macro}
% \end{macro}
%
% \begin{environment}{polycvletter}
% Here the actual letter environment is defined. 
% Like before, this is intended to be only one page and hence may cause problems for sufficiently long letters.
% The environment requires two parameters, the recipient address, and the body.
% First the page setup is changed, then the templates for header, recipient address, subject, and footer are called.
% The indentation is set to the specified value.
% At the end, the signature will be inserted, and below that - if present - the postscript remark.
%
%    \begin{macrocode}
\newenvironment{polycvletter}[3][\insertaddress]{%
  \newgeometry{%
    top=\polycvtopmargin,%
    left=2\polycvmargins,%
    right=2\polycvmargins,%
    bottom=\polycvbottommargin,%
    nohead,nofoot}%
  \polycv@letter@from{\insertauthor}{#1}%
  \polycv@letter@to{#3}%
  \polycv@letter@date%
  \polycv@letter@subject{#2}%
  \setlength{\parindent}{\polycvletterindent}%
}{%
  \vspace{0.5\baselineskip}%
  \noindent\includegraphics[width=5cm]{\insertsigfilename}\\%
  \noindent\insertauthor\par%
  \if\relax\detokenize{\polycv@letter@ps}\relax%
  \else%
  \polycv@letter@ps\par%
  \fi%
  \polycv@footer{}%
  \restoregeometry%
}
%    \end{macrocode}
% \end{environment}
%
% \subsection{Appendix}
%
% \begin{macro}{\pdfappendix}
% The appendix command is an interface to include (parts of) a pdf, e.g. for certificates or recommendations.
% The page colour is first reset to having none, otherwise it would overlay on top off the included pdf.
% By default all pages will be included, but this can be overwritten to  specific range (and other options).
% The command can be used multiple times.
%
%    \begin{macrocode}
\providecommand*\pdfappendix{}
\renewcommand*{\pdfappendix}[2][pages={1-}]{%
  \nopagecolor%
  \includepdf[#1]{#2}}%
%    \end{macrocode}
% \end{macro}
%
% \section{Building \pkg{polycv} and its Documentation}
% \label{sec:building}
%
% To compile the \pkg{polycv} package:
% \begin{enumerate}
%   \item \lstinline`pdflatex polycv.ins`
% \end{enumerate}
%
% To compile the \pkg{polycv} documentation
% \begin{enumerate}
%   \item \lstinline`pdflatex polycv.dtx` 
%   \item \lstinline`makeindex -s gglo.ist -o polycv.gls polycv.glo`
%   \item \lstinline`makeindex -s gind.ist -o polycv.ind polycv.idx`
%   \item \lstinline`pdflatex polycv.dtx` (several times)
% \end{enumerate}
% \clearpage
%
% \Finale
\endinput
